\documentclass[10pt,a4paper]{article}
\usepackage{amsmath,amssymb,amsthm,graphicx}
\usepackage{xfrac}
% Minimal editorial layout adapter for the pinned JEP source
% sources/symmetric-iet-infinite-jep-2025-302/source0.tex.
% The publisher document class (cedram) is not part of the retained material, and it is the
% class that used to define a handful of plain text abbreviations appearing in the author's
% prose. This adapter supplies exactly those abbreviations and nothing else: no mathematics,
% no notation, no theorem environment, no counter. The author's own macro block is retained
% separately from the pinned source itself.
% Corpus retention convention: CC BY 4.0, attribution to the original authors.
\providecommand{\ie}{i.e.\ }
\providecommand{\eg}{e.g.\ }
\providecommand{\resp}{resp.\ }
\providecommand{\skpt}{\smallskip}
\newenvironment{enumeratei}
{\bgroup\def\theenumi{\roman{enumi}}\def\theenumii{\arabic{enumii}}\begin{enumerate}}
{\end{enumerate}\egroup}

\hyphenation{addi-tion-al belong belongs denote exchan-ge exchan-ged exchan-ges}

\newcommand{\psfrac}[2]{\sfrac{(#1)}{#2}}

\newcommand\mto{\mathchoice{\longmapsto}{\mapsto}{\mapsto}{\mapsto}}
\let\oldbigsqcup\bigsqcup
\renewcommand{\bigsqcup}{\mathop{\textstyle\oldbigsqcup}\displaylimits}

\let\epsilon\varepsilon

\DeclareMathOperator{\Ess}{Ess}
\DeclareMathOperator{\Leb}{Leb}

\usepackage{tikz}
\usetikzlibrary{plotmarks}

\theoremstyle{plain}
\newtheorem{theorem}{Theorem}[section]
\newtheorem{lemma}[theorem]{Lemma}
\newtheorem{proposition}[theorem]{Proposition}
\newtheorem{corollary}[theorem]{Corollary}
\newtheorem{claim}{Claim}

\theoremstyle{definition}
\newtheorem{remark}[theorem]{Remark}
\newtheorem{example}[theorem]{Example}

\newcommand{\A}{\mathcal{A}}
\newcommand{\B}{\mathcal{B}}
\newcommand{\I}{\mathcal{I}}

\newcommand{\Z}{\mathbb{Z}}
\newcommand{\N}{\mathbb{N}}
\newcommand{\R}{\mathbb{R}}
\title{An ergodic symmetric interval exchange transformation whose middle point lies in a connection has minus one as an eigenvalue}
\author{Przemyslaw Berk, Frank Trujillo}
\date{doi:10.5802/jep.302; CC BY4.0}
\begin{document}
\maketitle
\noindent\textit{Editorial note.} The counted statement is the author's labelled corollary environment at source0.tex lines 893-895; its proof environment opens at line 897, the line after the blank separator, so the pairing is by adjacency and the retained proof range 897-902 is the whole environment.\par\medskip
\noindent\textit{Editorial note.} Four environments are retained as context because the counted proof depends on them: eq:backward\_return (source lines 565-568), prop:induced\_symmetric (source lines 573-579), prop:induced\_centers (source lines 583-590) and lem:sym\_int\_disjoint (source lines 728-731).\par\medskip
\noindent\textit{Editorial note.} This source has no prepared adapter and its preamble is built from the author's own macro and theorem block (source lines 8-45).\par\medskip
\noindent\textit{Editorial note.} This source has no label-number dump and the retained counter values are left empty, so the printed numbering of the card is the wrapper's default; the identity of the counted item is its author label.\par\medskip
\noindent\textit{Editorial note.} r002 note: r001 failed because the retained text uses the publisher class's plain text abbreviation resp (source line 901 in the proof, and source lines 575-576 in the retained prop:induced\_symmetric context), which the standalone wrapper does not define. r002 adds the minimal editorial adapter of this source as a first preamble span; the adapter contains only plain text abbreviations, no mathematics, no theorem environment and no counter, and nothing else in the card changes.\par\medskip
\noindent\textit{Editorial note.} r003 note: the r002 adapter span asked for line 13 of a twelve-line adapter file and the assembler refused it before compiling; r003 corrects the span to lines 1-12 and changes nothing else.\par\medskip
\noindent\textit{Editorial note.} r004 note: r003 failed with undefined references raised from inside the retained contexts themselves; r004 adds exactly the three environments those contexts refer to (cor: existenceofinfiniteorbit, eq: standardconjugacy and eq:delta\_1/2) and changes nothing else.\par\medskip
\noindent\textit{Editorial scope.} One original result of the pinned publisher article, counted once: minus one is an eigenvalue for an ergodic symmetric interval exchange transformation whose middle point lies in a connection (author label cor: evalue-1).\par\medskip
\noindent\textit{Allowed standard prerequisites (editorial).}\begin{itemize}
\item Koopman operator. The Koopman operator acts on functions by composition with the transformation, and its eigenvalues are the constants for which an eigenfunction exists.. It is the operator whose eigenvalue minus one is asserted.
\item Induced transformation on a symmetric subinterval. For a suitable symmetric subinterval the induced transformation is again an ergodic symmetric interval exchange transformation, and the centres of its exchanged intervals are images of the centres of the original one.. They provide the parity structure used to build the eigenfunction.
\item Rokhlin towers. The induced transformation decomposes the interval into towers of pairwise disjoint intervals of equal height.. The eigenfunction is defined level by level on those towers.
\end{itemize}\medskip
\begin{equation}
\label{eq:backward_return}
b_J(x) := \min\{m \geq 1 \mid T^{-m}(x) \in J\}.
\end{equation}

\begin{proposition}
\label{prop:induced_symmetric}
Let $T: I \to I$ be an ergodic symmetric IET and fix $\alpha \in \A \setminus \{\pi_0^{-1}(1)\}$ such that $M(\alpha) = +\infty$ (see Corollary \ref{cor: existenceofinfiniteorbit}). Fix $m \geq 1$ and let $J \subseteq I$ be the left-closed right-open subinterval with endpoints $T^{-m}(\partial I_\alpha)$, $T^{m}(\partial I_{\hat \alpha})$ (\resp $T^{-m + 1}(\partial I_\alpha)$, $T^{m}(\partial I_{\hat \alpha})$), where $\pi_0(\hat \alpha) = \pi_0(\alpha) - 1$.
Then $J$ is $\beta$-symmetric for some $\beta \in \A$ (\resp $\sfrac{1}{2}$-symmetric) and $T_J$ satisfies \eqref{eq: standardconjugacy}.

If, in addition, $J$ does not contain points from any connection. Then $T_J$ is an ergodic symmetric IET on $d - d'$ intervals, where $d'$ is the number of non-trivial {irreducible} connections of $T$.
\end{proposition}

\begin{proposition}
\label{prop:induced_centers}
Let $T$, $\alpha, \hat \alpha, \beta$ and $J$ as in Proposition \ref{prop:induced_symmetric}. Assume that $J$ does not contain points from any connection. Then the following holds:
\begin{enumerate}
\item For any $\gamma \in \A_J$ there exists unique $\sigma \in \A \cup \{\sfrac{1}{2}\}$, with $ \sigma \neq \beta$, and $\ell \geq 1$ such that $c_\gamma^J = p_J(c_\sigma) = T^{-\ell}(c_\sigma)$ and $T_J(p_J(c_\sigma)) = T^{\ell - \delta_{1/2}(\sigma)}(c_\sigma)$, where $\delta_{1/2}$ is given by \eqref{eq:delta_1/2}. Moreover, $c_\sigma$ does not belong to any non-trivial connection.
\item For any $\delta \in \A \cup \{\sfrac{1}{2}\}$ with $\delta \neq \beta$, either $p_J(c_\delta) = c_\gamma^J$ or $p_J(c_\delta) = \partial I_\gamma^J$ for some $\gamma \in \A_J$. In the latter case, $c_\delta$ lies inside a connection of $T$.
\end{enumerate}
\end{proposition}

\begin{lemma}
\label{lem:sym_int_disjoint}
Let $T: I \to I$ be an ergodic symmetric IET. Then there exists $\beta \in \A$ such that $c_\beta$ is not part of any non-trivial connection from $T$ and, for any $\epsilon > 0$, there exists a $\beta$-symmetric subinterval $J\subseteq I$ disjoint from the connections of $T$ satisfying $|J| < \epsilon$.
\end{lemma}

\begin{corollary}\label{cor: existenceofinfiniteorbit}
Assume that $T$ is an ergodic IET. Then there exists {$\beta \in\mathcal A \setminus \{\pi_0^{-1}(1)\}$ such that $M(\beta) = +\infty$. }
\end{corollary}

\begin{equation}\label{eq: standardconjugacy}
\mathcal{I}_I\circ T(x)=T^{-1}\circ\mathcal{I}_I(x) \quad \text{if } x \neq \partial I_{\alpha} \text{ for any } \alpha \in \A.
\end{equation}

\begin{equation}
\label{eq:delta_1/2}
\delta_{1/2}(\beta) = \left\{\begin{array}{ll} 1 & \text{ if } \quad \beta = \sfrac{1}{2}, \\
0 & \text{ otherwise}.
\end{array} \right.
\end{equation}

\begin{corollary}\label{cor: evalue-1}
Assume that $T$ is an ergodic symmetric IET such that $c_{1/2}$ lies inside a connection. Then $-1$ is an eigenvalue for the Koopman operator associated with $T$. In particular, $T$ is not weak mixing.
\end{corollary}

\begin{proof}
Let $J \subseteq I$ as in Proposition \ref{prop:induced_symmetric} such that $J$ does not contain any point from any connection. Such an interval exists by Lemma \ref{lem:sym_int_disjoint}. Then, by~Propo\-si\-tions~\ref{prop:induced_symmetric} and \ref{prop:induced_centers}, $T_J$ is symmetric and the middle points $\{c_\gamma^J\}_{\gamma \in \A_J}$ of the exchanged intervals $\{I_\gamma^J\}_{\gamma \in \A_J}$ are preimages of the middle points of the intervals exchanged by~$T_J$.
Moreover, again by Proposition \ref{prop:induced_centers}, the Rokhlin towers associated to $T_J$ are all of even height since for any $\gamma \in \A_J$ there exists $\sigma \in \A$ such that $c_\gamma^J = p_J(c_\sigma)$ and $T_J(c_\gamma^J) = T^{2b_J(c_\sigma)}(c_\gamma^J)$, where $b_J$ is given by \eqref{eq:backward_return}.

By defining a function $f$ that equals $1$ (\resp $-1$) on the odd (\resp even) levels of each Rokhlin tower, we~get an eigenfunction of $T$ with eigenvalue $-1$, that is, such that $f \circ T = -f$.
\end{proof}

\end{document}
